\documentclass[12pt,a4paper]{article}
\usepackage{ugmath}
\usepackage{float}
\usepackage{placeins}
\usepackage[labelfont=bf,labelsep=space]{caption}
\newcommand{\paren}[1]{\left( #1 \right)}
\newcommand{\alumno}{Ricardo León Martínez}
\newcommand{\materia}{Fisica I}
\newcommand{\profesor}{Lucero Uscanga Aguilera}
\newcommand{\tarea}{Tarea 5}
\newcommand{\fecha}{13/3/2026}

\begin{document}
    \begin{center}
    {\large\textbf{UNIVERSIDAD DE GUANAJUATO}}\\[0.3cm]
    {\normalsize\textbf{DIVISIÓN DE CIENCIAS NATURALES Y EXACTAS}}\\
    {\normalsize\textbf{CAMPUS GUANAJUATO}}\\[1cm]

    {\Large\textbf{\tarea\ (\materia)}}\\[1cm]
    \end{center}

    \textbf{Nombre:} \alumno \hfill 
    \textbf{Fecha:} \fecha \hfill 
    \textbf{Calificación:} \rule{3cm}{0.4pt} \\[0.3cm]

    \begin{problema}
        Una masa de 10kg se mueve bajo la acción de la fuerza $\vec{F}=\hat{u}_{x}(5t)+\hat{u}_{y}(3t^{2}-1)$N.
        Cuando $t=0$ el cuerpo está en reposo en el origen.
        \begin{enumerate}
            \item[(a)] Hallar el momento y la energía cinética del cuerpo cuando $t=10$s.
            \item[(b)] Calcular el impulso y el trabajo efectuado por la fuerza de $t=0$ a $t=10$s.
            Comparar con las repuestas en (a). 
        \end{enumerate}
    \end{problema}

    \begin{solucion}
        La fuerza aplicada es
        \begin{equation*}
            \vec{F}=\hat{u}_{x}(5t)+\hat{u}_{y}(3t^{2}-1).
        \end{equation*}
        La masa es
        \begin{equation*}
            m=10\text{ kg}.
        \end{equation*}
        Por la segunda ley de Newton,
        \begin{equation*}
            \vec{F}=m\vec{a}.
        \end{equation*}
        Luego
        \begin{equation*}
            \vec{a}(t)=\frac{\vec{F}(t)}{10}=\hat{u}_{x}\paren{\frac{t}{2}}+\hat{u}_{y}\paren{\frac{3t^{2}-1}{10}}.
        \end{equation*}
        Integrando la aceleración obtenemos la velocidad. Componente $x$
        \begin{equation*}
            v_{x}(t)=\int\frac{t}{2}dt=\frac{t^{2}}{4}+C_{1}.
        \end{equation*}
        Como el cuerpo parte del reposo,
        \begin{equation*}
            v_{x}(0)=0\implies C_{1}=0.
        \end{equation*}
        Así,
        \begin{equation*}
            v_{x}(t)=\frac{t^{2}}{4}.
        \end{equation*}
        Componente $y$
        \begin{equation*}
            v_{y}(t)=\int\frac{3t^{2}-1}{10}dt=\frac{t^{3}}{10}-\frac{t}{10}+C_{2}.
        \end{equation*}
        Usando $v_{y}(0)=0$,
        \begin{equation*}
            C_{2}=0.
        \end{equation*}
        Entonces
        \begin{equation*}
            v_{y}(t)=\frac{t^{3}-t}{10}.
        \end{equation*}
        Por tanto,
        \begin{equation*}
            \vec{v}(t)=\hat{u}_{x}\frac{t^{2}}{4}+\hat{u}_{y}\frac{t^{3}-t}{10}
        \end{equation*}
        \begin{enumerate}
            \item[(a)] Hallar el momento y la energía cinética del cuerpo cuando $t=0$s.\\
            Recordemos que
            \begin{equation*}
                \vec{p}=m\vec{v}.
            \end{equation*}
            En $t=10$
            \begin{gather*}
                v_{x}(10)=\frac{10^{2}}{4}=25,\\
                v_{y}(10)=\frac{10^{3}-10}{10}=\frac{990}{10}=99.
            \end{gather*}
            Así, 
            \begin{equation*}
                \vec{v}(10)=25\hat{u}_{x}+99\hat{u}_{y}.
            \end{equation*}
            Luego
            \begin{equation*}
                \vec{p}(10)=10(25\hat{u}_{x}+99\hat{u}_{y})=250\hat{u}_{x}+990\hat{u}_{y}\text{ kgm/s}.
            \end{equation*}
            Ahora la energia cinetica es
            \begin{equation*}
                K=\frac{1}{2}mv^{2}.
            \end{equation*}
            Primero calculamos
            \begin{equation*}
                v^{2}=25^{2}+99^{2}=625+9801=10426.
            \end{equation*}
            Entonces
            \begin{equation*}
                K=\frac{1}{2}(10)(10426)=52130\text{ J}.
            \end{equation*}
            Por tanto
            \begin{equation*}
                K(10)=52130\text{ J}
            \end{equation*}
            \item[(b)] Calcular el impulso y el trabajo efectuado por la fuerza de $t=0$ a $t=10$s.\\
            El impulso es
            \begin{gather*}
                \vec{\text{J}}_{x}=\int_{0}^{10}5tdt=\frac{5}{2}t^{2}\bigg|_{0}^{10}=250.\\
                \vec{\text{J}}_{y}=\int_{0}^{10}(3t^{2}-1)dt=(t^{3}-t)_{0}^{10}=1000-10=990.
            \end{gather*}
            Así,
            \begin{equation*}
                \vec{\text{J}}=250\hat{u}_{x}+990\hat{u}_{y}\text{ Ns}.
            \end{equation*}
            Como $1\text{ Ns}=1\text{ kgm/s}$,
            \begin{equation*}
                \vec{\text{J}}=\Delta\vec{p}
            \end{equation*}
            y esto coincide exactamente con el momento lineal obtenido en (a), pues el cuerpo
            inició en reposo. El trabajo total satisface el teorema trabajo-energía
            \begin{equation*}
                W=\Delta K.
            \end{equation*}
            Como incialmente el cuerpo estaba en reposo,
            \begin{equation*}
                K(0)=0.
            \end{equation*}
            Entonces
            \begin{equation*}
                W=K(10)-K(0)=52130.
            \end{equation*}
            Por tanto,
            \begin{equation*}
                W=52130\text{ J}
            \end{equation*}
            lo cual coincide exactamente con la energía cinética calculada en (a).
        \end{enumerate}
    \end{solucion}

    \begin{problema}
        Se lanza verticalmente hacia arriba un cuerpo de 20kg de masa con una velocidad de
        50m/s. Calcular
        \begin{enumerate}
            \item[(a)] Los valores iniciales de $E_{k},E_{p}$ y $E$
            \item[(b)] $E_{k}$ y $E_{p}$ despues de 3s
            \item[(c)] $E_{k}$ y $E_{p}$ a 100m de altura
            \item[(d)] La altura del cuerpo cuando $E_{k}$ es reducida un 80\% de su valor inicial.   
        \end{enumerate}
    \end{problema}

    \begin{solucion}
        Supondremos que el nivel de lanzamiento corresponde a energía potencial nula
        \begin{equation*}
            E_{p}(0)=0.
        \end{equation*}
        La energía mecánica es
        \begin{equation*}
            E=E_{k}+E_{p}.
        \end{equation*}
        Además,
        \begin{equation*}
            E_{k}=\frac{1}{2}mv^{2},\qquad E_{p}=mgh.
        \end{equation*}
        \begin{enumerate}
            \item[(a)] Los valores iniciales de $E_{k},E_{p}$ y $E$
            Inicialmente
            \begin{equation*}
                h_{0}=0,\qquad v_{0}=50\text{ m/s}.
            \end{equation*}
            Entonces
            \begin{equation*}
                E_{k}(0)=\frac{1}{2}(20)(50)^{2}.
            \end{equation*}
            Calculando
            \begin{equation*}
                E_{k}(0)=10(2500)=25000\text{ J}.
            \end{equation*}
            La energia potencial inicial es
            \begin{equation*}
                E_{k}(0)=20(9.8)(0)=0.
            \end{equation*}
            Por tanto, la energía total incial es
            \begin{equation*}
                E=25000\text{ J}.
            \end{equation*}
            Así,
            \begin{gather*}
                E_{k}(0)=25000\text{ J},\\
                E_{p}(0)=0,\\
                E=25000\text{ J}.
            \end{gather*}
            \item[(b)] $E_{k}$ y $E_{p}$ despues de 3s\\
            La velocidad en función del tiempo es
            \begin{equation*}
                v(t)=v_{0}-gt.
            \end{equation*}
            Entonces
            \begin{equation*}
                v(3)=50-9.8(3)=20.6\text{ m/s}.
            \end{equation*}
            La energía cinetica es
            \begin{gather*}
                E_{k}=\frac{1}{2}(20)(20.6)^{2}.\\
                E_{k}=10(424.36)=4243.6\text{ J}.
            \end{gather*}
            La altura tras 3 s es
            \begin{equation*}
                h(t)=v_{0}t-\frac{1}{2}gt^{2}.
            \end{equation*}
            Luego
            \begin{gather*}
                h(3)=50(3)-\frac{1}{2}(9.8)(9).\\
                h(3)=150-44.1=105.9\text{ m}.
            \end{gather*}
            Entonces
            \begin{gather*}
                E_{p}=20(9.8)(105.9).\\
                E_{p}=20756.4\text{ J}.
            \end{gather*}
            y efectivamente
            \begin{equation*}
                E_{k}+E_{p}=25000\text{ J}.
            \end{equation*}
            \item[(c)] $E_{k}$ y $E_{p}$ a 100m de altura\\
            La energía potencial a esa altura es
            \begin{equation*}
                E_{p}=20(9.8)(100).
            \end{equation*}
            Así,
            \begin{equation*}
                E_{p}=19600\text{ J}.
            \end{equation*}
            Como la energía mecánica se conserva,
            \begin{equation*}
                E_{k}=25000-19600.
            \end{equation*}
            Entonces
            \begin{equation*}
                E_{k}=5400\text{ J}.
            \end{equation*}
            Por tanto
            \begin{gather*}
                E_{p}=19600\text{ J}.\\
                E_{k}=5400\text{ J}.
            \end{gather*}
            \item[(d)] La altura del cuerpo cuando $E_{k}$ es reducida un 80\% de su valor inicial.\\
            La energía cinética inicial es
            \begin{equation*}
                E_{k0}=25000\text{ J}.
            \end{equation*}
            Reducirse un 80\% significa que queda el 20\%
            \begin{equation*}
                E_{k}=0.2(25000)=5000\text{ J}.
            \end{equation*}
            Por conservación de la energía,
            \begin{equation*}
                E_{p}=25000-5000=20000\text{ J}.
            \end{equation*}
            Entonces
            \begin{equation*}
                mgh=20000.
            \end{equation*}
            Sustituyendo
            \begin{gather*}
                20(9.8)h=20000.\\
                196h=20000.\\
                h=\frac{20000}{196}\approx102.04.
            \end{gather*}
            Así,
            \begin{equation*}
                h\approx102.04\text{ m}.
            \end{equation*}
        \end{enumerate}
    \end{solucion}

\end{document}